\section{Proposed Method}
\label{sec:methodology}

\subsection{Overview}

WR-STV is an STV-based anomaly scoring framework that transforms raw
service-path latency features into reliability-weighted residual
contributions. The method computes an anomaly score for each trace in
five stages, all derived from statistics learned from the training split
only: (1) construct the baseline Service Trace Vector (STV), (2) compute
per-feature training statistics, (3) convert observed values into
clipped standardized residuals, (4) weight each residual by a reliability
term reflecting how often that feature was observed during training, and
(5) aggregate the top-\(k\) weighted residuals into a single trace-level
score. 

Given a distributed trace, WR-STV produces a single anomaly score that
quantifies the degree to which the trace deviates from the statistical
behavior learned from training traces. Rather than modifying the underlying trace representation, WR-STV retains the original STV representation and instead modifies the anomaly scoring
process through feature-wise residual standardization, reliability
weighting, and top-\(k\) aggregation. This design preserves the interpretability of STV
features while reducing the influence of trace structural complexity on
the final anomaly score.


\subsection{Service Trace Vector Construction}

Following the Service Trace Vector (STV) representation introduced by
TraceAnomaly~\cite{liu2020traceanomaly}, each trace is mapped to a
fixed-vocabulary vector, where feature
\(j=(\text{service},\text{service\_path})\) is valued by
\(x_{i,j}=\log(1+d_{i,j})\). A sentinel value of \(-1\) is used when
feature \(j\) is not observed in trace \(i\). The feature vocabulary and
all statistics described in the following subsections are learned
exclusively from the training split. During inference, the same
vocabulary is reused without modification, and service paths not observed
during training are treated as absent rather than added as new features.
Using a fixed vocabulary ensures that training and inference operate over
a consistent feature space.

\subsection{Residual Normalization}

Raw STV values are not directly comparable across service-path features
because different paths naturally exhibit different latency scales and
variability. WR-STV therefore normalizes each observed feature relative
to its own training-time behavior. For each feature \(j\), we compute
the training-time mean and standard deviation over the training traces in
which \(j\) is observed:

\begin{align}
\mu_j
&=
\frac{1}{n_j}
\sum_{i \in \mathcal{T}_j}
x_{i,j},
\label{eq:mean}
\\
\sigma_j
&=
\sqrt{
\frac{1}{n_j}
\sum_{i \in \mathcal{T}_j}
(x_{i,j} - \mu_j)^2
},
\label{eq:std}
\end{align}

\noindent where \(\mathcal{T}_j\) is the set of training traces in which
feature \(j\) is present and \(n_j = |\mathcal{T}_j|\). If a feature has
insufficient observations or a numerically small standard deviation, the
implementation uses a unit standard deviation fallback to avoid unstable
division.

For an observed value \(x_{i,j}\) at inference time, the clipped
standardized residual is

\begin{equation}
\tilde{r}_{i,j}
=
\operatorname{clip}
\left(
\frac{x_{i,j}-\mu_j}{\sigma_j},
-C,
C
\right),
\label{eq:residual}
\end{equation}

\noindent where \(C\) denotes the residual clipping threshold. The
clipping operation improves numerical stability for features with very
small standard deviations and prevents individual feature deviations from
disproportionately influencing the final anomaly score. Features absent
from a trace are assigned zero residual contribution and therefore do not
increase the anomaly score.

\subsection{Reliability Weighting}

Features observed infrequently during training provide fewer samples for
estimating their underlying latency distribution, resulting in greater
uncertainty in the estimated mean and standard deviation. Consequently,
standardized residuals associated with such features are more susceptible
to sampling variability. To reduce the influence of these uncertain
estimates, each feature is assigned a reliability weight based on its
training-time observation frequency:

\begin{equation}
w_j
=
\frac{\log(1+n_j)}
{\max_{j'} \log(1+n_{j'})}.
\label{eq:weight}
\end{equation}

\noindent The logarithmic transformation compresses large count
differences while preserving the ordering of feature support. Thus,
\(w_j \in (0,1]\), with the most frequently observed feature receiving
weight \(1\). Features observed more frequently during training receive
larger weights, whereas infrequently observed features are down-weighted,
reducing their influence on the final anomaly score.

\subsection{Top-\(k\) Weighted Residual Scoring}

The weighted contribution of feature \(j\) in trace \(i\) is defined as

\begin{equation}
c_{i,j}
=
w_j
\left|
\tilde{r}_{i,j}
\right|.
\label{eq:contribution}
\end{equation}

\noindent Larger standardized deviations from learned normal behavior
produce larger contributions, moderated by the feature reliability
weight.

Rather than summing all feature contributions, which can reintroduce
dependence on the number of populated features in a trace, WR-STV
aggregates only the strongest residual evidence. The final anomaly score
is computed as the average of the \(k\) largest weighted contributions:

\begin{equation}
\operatorname{score}(i)
=
\frac{1}{k}
\sum_{j \in \operatorname{Top}_k(c_{i,:})}
c_{i,j}.
\label{eq:topk}
\end{equation}

\noindent This aggregation emphasizes the most significant service-path
deviations while limiting the influence of numerous low-magnitude
residuals that are unlikely to indicate anomalous behavior. The value of
\(k\) is selected through validation, as described in
Section~\ref{sec:hyperparameter-selection}. In the final configuration,
we use \(k=1\). Under this setting, the anomaly score corresponds to the
largest weighted residual within the trace, and the associated
service-path feature is reported as the explanation for the score in the
interpretability analysis presented in
Section~\ref{sec:results-explainability}.

\subsection{Request-Aware Calibration (Ablation Variant)}

As an ablation study, we also evaluate a request-aware variant in which
the statistics in Equations~\ref{eq:mean}--\ref{eq:std} are computed
separately for each root operation when sufficient training samples are
available. When insufficient data exist for a particular request type,
the corresponding global statistics are used instead.

Although this variant further reduces the observed correlation between
anomaly score and trace complexity
\begin{table}[t]
\centering
\caption{Complexity-bias correlation across methods, averaged over
num\_services and num\_spans (real, non-injected test traces). This is the
paper's central bias-reduction evidence.}
\label{tab:complexity-bias-all-methods}
\begin{tabular}{lcc}
\toprule
\textbf{Method} & \textbf{corr.\ num\_services} & \textbf{corr.\ num\_spans} \\
\midrule
Baseline STV + IF                 & 0.868 & 0.820 \\
Enhanced STV + IF                 & 0.835 & 0.779 \\
Global Residual STV (unweighted)  & 0.383 & 0.324 \\
\textbf{WR-STV (weighted)}        & \textbf{0.296} & \textbf{0.238} \\
Request-Aware WR-STV              & 0.274 & 0.218 \\
\bottomrule
\end{tabular}
\end{table}

(Table~\ref{tab:complexity-bias-all-methods}), it does not consistently improve
anomaly scoring performance. Therefore, we report it as an ablation
rather than the primary WR-STV method.
